\begin{claim}{F.15}{Proved}{The two-point kernel and its full involution cocycles}
Retain the actual based configuration of F.12--14 and put
$c_0=v_2(s_0)$, $x_0=v_0=1$, $y_0=c_0$.  On the sphere
with the other six labels retained, set
\[
\begin{split}
 M&=\mathbf P^1\setminus
       \{0,\infty,v_1(s_0),v_3(s_0),v_4(s_0),v_5(s_0)\},\\
 r(z)&=c_0/z,\qquad
 X=\operatorname{Conf}_2(M),\qquad
 \mathcal R(x,y)=(r(y),r(x)).
\end{split}                                                     \tag{F.15.1}
\]
Thus $\mathcal R(x_0,y_0)=(x_0,y_0)$.  Write $a_+,a_-$
for the two roots of $z^2=c_0$, which are unmarked points of
$M$.  Write $\mathsf r$ for the actual eight-puncture
mapping class of $r$, so conjugation by $\mathsf r$ induces
the based automorphism $R_*$ of F.12.

There are the following actual completed homomorphisms:
\[
\begin{split}
 J&=\pi_1(X,(x_0,y_0)),\\
 L_0&=\pi_1(M,x_0),\qquad L_2=\pi_1(M,y_0),\\
 1\longrightarrow B\longrightarrow\widehat J
   \xrightarrow{\ p\ }\widehat L_0\times\widehat L_2
     \longrightarrow1 .
\end{split}                                                     \tag{F.15.2}
\]
Here $\widehat J$ is the kernel in
$P_8=\widehat{P\Gamma}_{0,8}$ of forgetting the pair
$\{v_0,v_2\}$; $p$ records the two labelled point tracks,
and its kernel is exactly $B=K_0\cap K_2$ of F.14.
Let $\rho=r_*:\widehat L_0\longrightarrow\widehat L_2$.
The induced involution and its fixed subgroup are
\[
 R_*(u,v)=(\rho^{-1}(v),\rho(u)),\qquad
 (\widehat L_0\times\widehat L_2)^{R_*}
        =\{(u,\rho(u)):u\in\widehat L_0\}.                       \tag{F.15.3}
\]

The actual fixed-configuration map
\[
 i:M\setminus\{a_+,a_-\}\longrightarrow X,\qquad
                         z\longmapsto(z,r(z))                 \tag{F.15.4}
\]
induces a full profinite embedding, and
\[
\begin{split}
 C_{\widehat J}(\mathsf r)
   &=\widehat i_*\,
        \widehat{\pi_1(M\setminus\{a_+,a_-\},x_0)}
       \ \cong\widehat F_7,\\
 p\bigl(C_{\widehat J}(\mathsf r)\bigr)
   &=(\widehat L_0\times\widehat L_2)^{R_*},\\
 C_B(\mathsf r)
   &\cong\ker\bigl(
       \widehat{\pi_1(M\setminus\{a_+,a_-\},x_0)}
           \xrightarrow{\ \mathrm{fill}_{a_+,a_-}\ }
                  \widehat{\pi_1(M,x_0)}\bigr).
\end{split}                                                     \tag{F.15.5}
\]
In particular the last kernel is the closed normal subgroup
generated by the two based branch meridians in this
specified $\widehat F_7$.

Every full continuous $R_*$-cocycle in $B$ is a
$B$-coboundary.  More precisely, there is a homeomorphism
of compact spaces
\[
\begin{split}
 C_B(\mathsf r)\backslash B
    &\xrightarrow{\ \sim\ }
        Z^1(R_*,B):=\{\delta\in B:\delta R_*(\delta)=1\},\\
 C_B(\mathsf r)b&\longmapsto b^{-1}R_*(b).
\end{split}                                                     \tag{F.15.6}
\]
Consequently the quotient map $\pi:P_8\to P_8/B$ gives the
exact sequence
\[
 1\longrightarrow C_B(\mathsf r)
  \longrightarrow C_{P_8}(\mathsf r)
  \xrightarrow{\ \pi\ }(P_8/B)^{R_*}
  \longrightarrow1 .                                          \tag{F.15.7}
\]

For the exact equalizer $\mathcal E_7$ of F.14, set
\[
 \mathcal B_u=\{b\in B:b^{-1}R_*(b)=d(u)\}.
                                                               \tag{F.15.8}
\]
This is a nonempty left $C_B(\mathsf r)$-coset for every
$u\in\mathcal E_7$.  The centralizing representatives of
its actual $B$-coset are exactly
\[
 \bigl(\widehat\phi(u)B\bigr)\cap C_{P_8}(\mathsf r)
       =\widehat\phi(u)b^{-1}C_B(\mathsf r)
                 \qquad(b\in\mathcal B_u).                    \tag{F.15.9}
\]
The source embedding has trivial intersection with $B$.
It follows that the intersection of (F.15.9) with
$\widehat\phi(\widehat{\pi_1(Y)})$ is either empty or the
singleton $\{\widehat\phi(u)\}$, and is the latter
exactly for $u\in\mathcal E$.  These formulas determine
the correction and its ambiguity; triviality of the
full residual group $\mathcal E$ is still unresolved.
\end{claim}

\begin{proof}
\emph{The labelled configuration and the two tracks.}
The actual quotient representative in F.12 is
$z\mapsto v_2(s_0)/z$.  Its permutation on the eight
labels is
\[
       (0\ \infty)(v_0\ v_2)(v_1\ v_3)(v_4\ v_5).
                                                               \tag{F.15.10}
\]
It has no fixed marked point.  Its two fixed points are
therefore precisely the unmarked $a_+,a_-$ of (F.15.1).
Deleting the invariant pair $\{v_0,v_2\}$ retains an
actual homeomorphism of the indicated six-puncture
sphere, with no additional base path.
Under point pushing, conjugation by its mapping class
acts on the ordered two-point fibre by $\mathcal R$.
This proves both the specified action and (F.15.3).

The discrete Birman fibration identifies $J$ with the
kernel of forgetting those two labels.  Projecting $X$
to one coordinate has fibre a sphere with seven
punctures, and gives a split extension
\[
              1\longrightarrow F_6
                \longrightarrow J
                \longrightarrow F_5
                \longrightarrow1 .                            \tag{F.15.11}
\]
One continuous section can be seen in a coordinate with
one retained puncture at infinity: put the second point
on the positive real axis beyond
$2+|x|+\sum_{z\in P_6\setminus\{\infty\}}|z|$.
A path in the fibre adjusts its basepoint to the
specified one.  The maps in (F.15.11), including that
adjustment, remain the ordinary labelled forgetful maps.

The finitely generated free kernel in (F.15.11) has its
full induced profinite topology.  Indeed, its
characteristic finite-index subgroups are cofinal among
its finite-index subgroups.  For any such subgroup, the
quotient of the free kernel is finite and the action of
the free quotient on it has finite image; the associated
finite semidirect product supplies the required finite
quotient of $J$.  Completing (F.15.11) is therefore
exact.
Map this completed sequence into the two successive
completed Birman sequences of F.14.14.  On its free
kernel and its free quotient the maps are the already
specified embeddings $\widehat F_6$ and $\widehat F_5$.
An element in the kernel of $\widehat J\to P_8$ must
first map trivially to $\widehat F_5$, then trivially to
$\widehat F_6$.  The map is injective.
Conversely, an element in the kernel of forgetting both
points can first be matched in the $\widehat F_5$
quotient by an element of $\widehat J$; the remaining
difference lies in the embedded $\widehat F_6$.
This proves the asserted equality with the actual
completed two-point kernel.

Each coordinate projection of $J$ records the track of
that named point.  The fibre of the first projection
maps onto $L_2$ by filling $x_0$, so the product track
map is onto $L_0\times L_2$.  The same is true after
completion.  The completion of a direct product is the
product of the full completions: any homomorphism to a
finite group restricts to homomorphisms on its two
commuting factors, and their continuous extensions
multiply to its unique continuous extension.
The kernel of the product track map is the simultaneous
kernel of the two point-forgetting restrictions.
Inside the embedded $\widehat F_6$ it is precisely
the kernel of filling the other named point.
The last Birman square in F.14 identifies it with
$K_0\cap K_2$, including its closed normal
meridian description.  This proves (F.15.2) without a
closure--intersection assumption.

\emph{The discrete fixed subgroup from the actual quotient cover.}
Consider the specified double cover obtained by deleting
the two fixed points of $r$:
\[
 M\setminus\{a_+,a_-\}\longrightarrow
   (M/\langle r\rangle)\setminus\{\overline a_+,
                                      \overline a_-\}.
                                                               \tag{F.15.12}
\]
The target is a sphere with three ordinary punctures
and two branch points removed.  Its group is $F_4$.
The monodromy takes the two branch meridians to the
nonidentity element of $C_2$ and takes the three
ordinary puncture meridians to the identity.
Its kernel is the actual $\pi_1$ of the source of
(F.15.12), of index two and rank seven.

Use the discrete covering-normalizer sequence F.7.5
for $r$ before and after forgetting the pair
$\{x_0,y_0\}$.  Downstairs this forgets its one ordinary
orbit point.  The kernel of this downstairs
forgetting map is the point-pushing $F_4$ just
specified.  A point-pushing isotopy lifts from the
identity through the double cover, fixing all retained
punctures individually.  Its endpoint fixes both
forgotten labels individually precisely when its
monodromy in (F.15.12) is trivial.
The finite deck kernel has trivial intersection with
the pure upstairs group, because $r$ exchanges the
retained punctures in three pairs.
Conversely, an element of $J$ that centralizes $r$
descends through this normalizer sequence.  After
forgetting the orbit point its descent is the identity,
since its upstairs filling is the identity and the
quotient map has only the displayed deck kernel.
It is therefore exactly one of these monodromy-zero
point pushes.  This proves
\[
 C_J(\mathsf r)
      =i_*\pi_1(M\setminus\{a_+,a_-\},x_0)
      \cong F_7.                                               \tag{F.15.13}
\]
The lift of its downstairs point track is the path
$(z,r(z))$.  Thus the equality uses the actual map
(F.15.4), not merely an abstract isomorphism of free
groups.  Projection to the first track is exactly
filling the two fixed points, and projection to the
second is its composite with $\rho$.

This $F_7$ also has its full induced profinite topology.
The downstairs point-pushing $F_4$ has its full topology
in the decorated quotient mapping class group by the
completed Birman argument of F.14.  Its index-two
monodromy kernel inherits its full topology.
The pure centralizer is a finite-index subgroup of the
normalizer of the specified order-two group, and its
descent is injective because the deck element is
nonpure.  Its image is finite index in the decorated
quotient, so this restriction is an isomorphism onto
that subgroup.  The prime-order normalizer topology in
F.7.4--5 consequently transports the full topology of
that $F_7$ to its actual subgroup in the eight-puncture
mapping class group.  It remains full in its subgroup
$J$ as well.  Hence the completion of (F.15.13)
is an embedding of this specified $\widehat F_7$ in
$\widehat J$.

\emph{One discrete conjugacy class of labelled sections.}
Let
\[
                       E=J\rtimes\langle\mathsf r\rangle
                                                               \tag{F.15.14}
\]
with the actual action already specified.
It embeds in the eight-puncture mapping class group:
the nontrivial coset is detected by the finite label
permutation (F.15.10).
Every order-two section in $E$ has that same permutation
and fills to the specified mapping class of $r$ on $M$.

We first justify pure conjugacy in the eight-puncture
group.  The cyclic case of Nielsen realization, in
the punctured-surface form of
\cite[Theorem 4 and Section IV, pp.~255--256]{bKerckhoff1983},
realizes such a section as an orientation-preserving
involution of the punctured sphere.  It extends over
the punctures.  With the invariant conformal structure,
filling gives a Riemann sphere, and a nontrivial
holomorphic involution has two fixed points.
Neither is marked, by (F.15.10).
Its quotient and that of $r$ are spheres with two
branch points and four specified ordinary orbit
points.  Choose an orientation-preserving quotient
homeomorphism matching those branch points and the
four named orbit pairs.  Both double covers have
monodromy one at the branch points and zero at the
ordinary points, so it lifts to an equivariant
homeomorphism.

The lift may interchange the two names inside some
of the four pairs.  Each mismatch can be corrected
independently: in the target quotient push that
ordinary orbit point once around one branch point,
holding the other orbit points and both branch points
fixed.  Lift the isotopy starting at the identity.
The endpoint commutes with $r$, interchanges exactly
the two labels in the selected orbit, and fixes the
other six labels individually.  The covering
monodromy proves these exact permutations.
Composing the required corrections yields a pure
conjugator.  Thus every section is $g\mathsf r g^{-1}$
with $g\in P\Gamma_{0,8}$.

The actual pure centralizer forgetting map is onto:
\[
 C_{P\Gamma_{0,8}}(\mathsf r)\longrightarrow
                              C_{P\Gamma_{0,6}}(r).
                                                               \tag{F.15.15}
\]
To prove this, take a class on the right.  The
discrete covering-normalizer map realizes it by a
homeomorphism commuting with the actual $r$.
Its value at $x_0$ is in the connected surface
$M\setminus\{a_+,a_-\}$.  Choose a path in that
surface from this value to $x_0$.  The quotient
path moves the one ordinary orbit point while
avoiding all branch points and retained punctures.
An isotopy moving that point, supported away from
those fixed labels, lifts from the identity to an
$r$-equivariant isotopy.  Its chosen lift carries
the image of $x_0$ to $x_0$ and the paired point to
$y_0$.  Composing its endpoint with the original
representative gives a pure eight-label centralizer
representative with the same six-label class.
This proves (F.15.15).

For a section $s=g\mathsf r g^{-1}$ in $E$, its
filled class is $r$, so the filled image of $g$
centralizes $r$.  Choose $c$ in the left side of
(F.15.15) with that same image.  Then $h=gc^{-1}$
belongs to $J$ and
$s=h\mathsf r h^{-1}$.  Consequently
\[
          \{\hbox{order-two sections of }E\}/J
                            \quad\hbox{has one element}.       \tag{F.15.16}
\]

\emph{The hypotheses for full profinite section comparison.}
We verify them for this particular $E$, rather than
assuming that goodness passes to an arbitrary
subgroup of a mapping class group.
The split sequence (F.15.11) and its exact completion
show that $J$ is residually finite: an element with
nontrivial image in $F_5$ is detected there, and an
element of the free kernel is detected by one of
the finite semidirect products used above.
The finitely generated group $J$ has cofinal
characteristic finite-index subgroups.  Taking
these to be invariant under the order-two
automorphism shows directly that the finite
extension $E$ is residually finite and that
$\widehat E=\widehat J\rtimes\langle\mathsf r\rangle$.

The free kernel and free quotient in (F.15.11)
are good.  With a finite coefficient module their
cohomology is finite and vanishes above degree one.
The discrete and profinite
Lyndon--Hochschild--Serre spectral sequences for
(F.15.11) therefore have isomorphic second pages:
the coefficients $H^b(F_6,A)$ for $b=0,1$ are
finite, so goodness of $F_5$ applies to them.
This proves goodness of $J$ and finiteness of its
cohomology with every finite module.
The same spectral-sequence comparison for
$1\to J\to E\to C_2\to1$ proves these assertions
for $E$.
The free-by-free group $J$ has a finite
two-dimensional classifying complex: take a
finite graph for the free kernel over a graph
for the free quotient, using the finitely many
monodromy homotopy equivalences on its edges.
Its universal cover is a tree of contractible
covering spaces.  In particular
$\operatorname{cd}(J)\leq2$.
The resulting finite-type resolutions and the
finite-group resolution for $C_2$ show that
cohomology of $E$ commutes with filtered unions
of coefficient modules.  A discrete
$\mathbf F_2[\![\widehat E]\!]$-module is the
union of its finite invariant submodules:
the span of each finite orbit is finite.
Thus the finite-coefficient comparison just
proved supplies the comparison for all
discrete modules required in
\cite[Theorem A]{R-BZ}.
The virtual cohomological dimension and
finite-cohomology hypotheses there are now
verified as well.

Apply that prime-order theorem to
$E\hookrightarrow\widehat E$.
It identifies the conjugacy classes of its
finite $2$-subgroups and gives the closure
of the centralizer of $\mathsf r$.
The group $\widehat J$ is torsion-free, by
its exact extension of the two torsion-free
free profinite groups in (F.15.11).
Hence every order-two subgroup of $\widehat E$
is a section.  By (F.15.16) there is one such
conjugacy class.  A conjugator in $\widehat E$
can be replaced by one in $\widehat J$:
write it as $h\mathsf r^\epsilon$ and use
that $\mathsf r^\epsilon$ centralizes
$\mathsf r$.
The centralizer comparison, intersected with
the open subgroup $\widehat J$, gives
\[
 C_{\widehat J}(\mathsf r)
                  =\overline{C_J(\mathsf r)}.
                                                               \tag{F.15.17}
\]
Combining this equality with the exact
full-topology statement following (F.15.13)
proves the first line of (F.15.5).
Its projection is precisely the graph of
the full filling epimorphism
$\widehat F_7\to\widehat F_5$ and its
composite with $\rho$.  This proves the
other two lines, including the named kernel.

\emph{Removing the fixed quotient coordinate.}
Take $\delta\in B$ with
$\delta R_*(\delta)=1$.
Then $\delta\mathsf r$ is an order-two
section of $\widehat E$.  The single
section class just proved gives $h\in\widehat J$
with
\[
                  \delta=h^{-1}R_*(h).
                                                               \tag{F.15.18}
\]
Its image $p(h)$ is fixed by (F.15.3).
By the second line of (F.15.5), choose
$c\in C_{\widehat J}(\mathsf r)$ with
$p(c)=p(h)$.  Then $b=c^{-1}h$ lies in
$B$ and direct multiplication yields
$b^{-1}R_*(b)=\delta$.
This proves surjectivity in (F.15.6).
Two such elements $b,b'$ have the same
coboundary exactly when
$R_*(b'b^{-1})=b'b^{-1}$, or equivalently
$b'\in C_B(\mathsf r)b$.  This proves the
asserted fibres.  The map is continuous,
its source is compact, and the cocycle
set is a closed subspace of the profinite
group $B$.  The bijection is therefore
a homeomorphism.

The subgroup $B$ is closed and normal in
$P_8$, and $R_*$ preserves it by F.14.13.
For any fixed element of $P_8/B$, choose a
lift $z\in P_8$.  Its defect
$z^{-1}R_*(z)$ is an $R_*$-cocycle in $B$.
Choose $b$ with the same coboundary by
(F.15.6).  Then
$R_*(zb^{-1})=zb^{-1}$, giving a
centralizing lift with the same quotient
image.  Its kernel is exactly
$C_B(\mathsf r)$.
This proves every map and the surjectivity
in (F.15.7).

\emph{The original source and its remaining membership equation.}
F.14 places $d(u)$ in the cocycle set
of (F.15.6), so (F.15.8) is nonempty
and has the specified left coset form.
For any one of its elements,
\[
 R_*(\widehat\phi(u)b^{-1})
  =\widehat\phi(u)d(u)R_*(b)^{-1}
  =\widehat\phi(u)b^{-1}.
                                                               \tag{F.15.19}
\]
Any two centralizing lifts of the same
$B$-coset differ on the right by an
element of $C_B(\mathsf r)$, proving
(F.15.9).
Finally $B\subset K_0$, and F.14 proves
that the composite of $\widehat\phi$
with forgetting $v_0$ is injective.
Thus $\widehat\phi(\widehat{\pi_1(Y)})\cap B=1$.
If two source elements have images in the
same $B$-coset, the same injectivity
makes them equal.
The only possible source representative
of (F.15.9) is consequently
$\widehat\phi(u)$ itself; it centralizes
$\mathsf r$ precisely for
$u\in\mathcal E$, by F.12.5.
No choice of the correcting element
$b$ has been asserted to remain in
the original source image.
\end{proof}
